%% style/environments.tex
%% The five locked tcolorbox environments + Colab box + warm-up + exercise.
%%
%% Discipline: do NOT add new colored environments without explicit consent.
%% Five is the design budget.
%%
%% Counter discipline: definitions, theorems, propositions, lemmas, corollaries,
%% and examples all SHARE one counter (`mathemagicsCounter`) which resets per
%% chapter. So you'll see "Definition 4.1, Theorem 4.2, Example 4.3" sequentially
%% — easy to find by number, no risk of collisions.

% Master shared counter
\newcounter{mathemagicsCounter}[chapter]
\renewcommand{\themathemagicsCounter}{\thechapter.\arabic{mathemagicsCounter}}

%%%%%%%%%%%%%%%%%%%%%%%%%%%%%%%%%%%%%%%%%%%%%%%%%%%%%%%%%%%%%%%%%%%%%%%%%%%%%%%
% (1) Definition — light blue tinted box, blue rule
%%%%%%%%%%%%%%%%%%%%%%%%%%%%%%%%%%%%%%%%%%%%%%%%%%%%%%%%%%%%%%%%%%%%%%%%%%%%%%%

\newtcbtheorem[use counter=mathemagicsCounter]
  {definition}{Definition}{%
    enhanced,breakable,
    colback=mathemagicsDefBg,
    colframe=mathemagicsDef,
    boxrule=0.6pt,
    arc=2pt,
    fonttitle=\bfseries,
    coltitle=mathemagicsDef,
    colbacktitle=mathemagicsDefBg,
    attach boxed title to top left={yshift=-2mm, xshift=4mm},
    boxed title style={colback=mathemagicsDefBg, colframe=mathemagicsDef,
                       boxrule=0.6pt, sharp corners=west, arc=2pt},
    top=4mm,
    before skip=1em, after skip=1em
  }{def}

%%%%%%%%%%%%%%%%%%%%%%%%%%%%%%%%%%%%%%%%%%%%%%%%%%%%%%%%%%%%%%%%%%%%%%%%%%%%%%%
% (2) Theorem / Proposition / Lemma — minimal, white box, slanted body, shared counter
%%%%%%%%%%%%%%%%%%%%%%%%%%%%%%%%%%%%%%%%%%%%%%%%%%%%%%%%%%%%%%%%%%%%%%%%%%%%%%%

\newtcbtheorem[use counter=mathemagicsCounter]
  {theorem}{Theorem}{%
    enhanced,breakable,
    colback=mathemagicsThmBg,
    colframe=mathemagicsThm,
    boxrule=0.8pt,
    arc=0pt,
    fonttitle=\bfseries\itshape,
    coltitle=mathemagicsThm,
    colbacktitle=mathemagicsThmBg,
    attach boxed title to top left={yshift=-2mm, xshift=4mm},
    boxed title style={colback=mathemagicsThmBg, colframe=mathemagicsThm,
                       boxrule=0.8pt, sharp corners, arc=0pt},
    fontupper=\itshape,
    top=4mm,
    before skip=1em, after skip=1em
  }{thm}

\newtcbtheorem[use counter=mathemagicsCounter]
  {proposition}{Proposition}{%
    enhanced,breakable,
    colback=mathemagicsThmBg,
    colframe=mathemagicsThm,
    boxrule=0.8pt,
    arc=0pt,
    fonttitle=\bfseries\itshape,
    coltitle=mathemagicsThm,
    colbacktitle=mathemagicsThmBg,
    attach boxed title to top left={yshift=-2mm, xshift=4mm},
    boxed title style={colback=mathemagicsThmBg, colframe=mathemagicsThm,
                       boxrule=0.8pt, sharp corners, arc=0pt},
    fontupper=\itshape,
    top=4mm,
    before skip=1em, after skip=1em
  }{prop}

\newtcbtheorem[use counter=mathemagicsCounter]
  {lemma}{Lemma}{%
    enhanced,breakable,
    colback=mathemagicsThmBg,
    colframe=mathemagicsThm,
    boxrule=0.8pt,
    arc=0pt,
    fonttitle=\bfseries\itshape,
    coltitle=mathemagicsThm,
    colbacktitle=mathemagicsThmBg,
    attach boxed title to top left={yshift=-2mm, xshift=4mm},
    boxed title style={colback=mathemagicsThmBg, colframe=mathemagicsThm,
                       boxrule=0.8pt, sharp corners, arc=0pt},
    fontupper=\itshape,
    top=4mm,
    before skip=1em, after skip=1em
  }{lem}

\newtcbtheorem[use counter=mathemagicsCounter]
  {corollary}{Corollary}{%
    enhanced,breakable,
    colback=mathemagicsThmBg,
    colframe=mathemagicsThm,
    boxrule=0.8pt,
    arc=0pt,
    fonttitle=\bfseries\itshape,
    coltitle=mathemagicsThm,
    colbacktitle=mathemagicsThmBg,
    attach boxed title to top left={yshift=-2mm, xshift=4mm},
    boxed title style={colback=mathemagicsThmBg, colframe=mathemagicsThm,
                       boxrule=0.8pt, sharp corners, arc=0pt},
    fontupper=\itshape,
    top=4mm,
    before skip=1em, after skip=1em
  }{cor}

%%%%%%%%%%%%%%%%%%%%%%%%%%%%%%%%%%%%%%%%%%%%%%%%%%%%%%%%%%%%%%%%%%%%%%%%%%%%%%%
% (3) Intuition — cream box, no number, used liberally to translate theorems
%%%%%%%%%%%%%%%%%%%%%%%%%%%%%%%%%%%%%%%%%%%%%%%%%%%%%%%%%%%%%%%%%%%%%%%%%%%%%%%

\newtcolorbox{intuition}{%
  enhanced,breakable,
  colback=mathemagicsIntuBg,
  colframe=mathemagicsIntu,
  boxrule=0.4pt,
  arc=2pt,
  leftrule=2.5pt,
  rightrule=0.4pt,
  toprule=0.4pt,
  bottomrule=0.4pt,
  fonttitle=\bfseries\sffamily\small,
  title={Intuition},
  attach boxed title to top left={yshift=-2mm, xshift=4mm},
  boxed title style={colback=mathemagicsIntu, colframe=mathemagicsIntu,
                     sharp corners, boxrule=0pt},
  coltitle=white,
  top=4mm,
  before skip=0.8em,
  after skip=0.8em,
}

%%%%%%%%%%%%%%%%%%%%%%%%%%%%%%%%%%%%%%%%%%%%%%%%%%%%%%%%%%%%%%%%%%%%%%%%%%%%%%%
% (4) Example — pale gray, shared counter
%%%%%%%%%%%%%%%%%%%%%%%%%%%%%%%%%%%%%%%%%%%%%%%%%%%%%%%%%%%%%%%%%%%%%%%%%%%%%%%

\newtcbtheorem[use counter=mathemagicsCounter]
  {example}{Example}{%
    enhanced,breakable,
    colback=mathemagicsExBg,
    colframe=mathemagicsEx,
    boxrule=0.4pt,
    arc=2pt,
    fonttitle=\bfseries\sffamily,
    coltitle=white,
    colbacktitle=mathemagicsEx,
    attach boxed title to top left={yshift=-2mm, xshift=4mm},
    boxed title style={colback=mathemagicsEx, colframe=mathemagicsEx,
                       sharp corners, boxrule=0pt},
    top=4mm,
    before skip=1em, after skip=1em
  }{ex}

%%%%%%%%%%%%%%%%%%%%%%%%%%%%%%%%%%%%%%%%%%%%%%%%%%%%%%%%%%%%%%%%%%%%%%%%%%%%%%%
% (5) Aside — uses tufte-book's \marginnote.
% Three flavors: \asideHistorical, \asideMatters, \asideCartoon.
%%%%%%%%%%%%%%%%%%%%%%%%%%%%%%%%%%%%%%%%%%%%%%%%%%%%%%%%%%%%%%%%%%%%%%%%%%%%%%%

\makeatletter
\newcommand{\@asideTitleStyle}[1]{%
  {\sffamily\bfseries\scshape\small #1}%
}
\makeatother

\makeatletter
\newcommand{\asideHistorical}[1]{%
  \marginnote{\@asideTitleStyle{Historical aside.} \footnotesize #1}%
}
\newcommand{\asideMatters}[1]{%
  \marginnote{\@asideTitleStyle{Why this matters.} \footnotesize #1}%
}
\newcommand{\asideCartoon}[1]{%
  \marginnote{\@asideTitleStyle{Aside.} \footnotesize #1}%
}
\newcommand{\aside}[2][Aside]{%
  \marginnote{\@asideTitleStyle{#1.} \footnotesize #2}%
}
\makeatother

%%%%%%%%%%%%%%%%%%%%%%%%%%%%%%%%%%%%%%%%%%%%%%%%%%%%%%%%%%%%%%%%%%%%%%%%%%%%%%%
% Code: three placements
%   1. \mintinline{python}|...|     (inline) — provided by minted package
%   2. \margincode{...}              (margin)
%   3. \begin{colabbox}...\end{colabbox}  (full-width experiment)
%%%%%%%%%%%%%%%%%%%%%%%%%%%%%%%%%%%%%%%%%%%%%%%%%%%%%%%%%%%%%%%%%%%%%%%%%%%%%%%

\newcommand{\margincode}[1]{%
  \marginnote{%
    \begin{tcolorbox}[%
      colback=mathemagicsCodeBg,
      boxrule=0pt,
      arc=1pt,
      left=2pt,right=2pt,top=2pt,bottom=2pt,
      fontupper=\footnotesize\ttfamily,
    ]
      #1
    \end{tcolorbox}%
  }%
}

\newtcolorbox{colabbox}[1][]{%
  enhanced,breakable,
  colback=mathemagicsCodeBg,
  colframe=mathemagicsColab,
  boxrule=0.4pt,
  arc=2pt,
  fonttitle=\bfseries\sffamily\small,
  title={Simulation},
  attach boxed title to top left={yshift=-2mm, xshift=4mm},
  boxed title style={colback=mathemagicsColab, colframe=mathemagicsColab,
                     sharp corners, boxrule=0pt},
  coltitle=black,
  top=4mm,
  before skip=1em, after skip=1em,
  #1
}

%%%%%%%%%%%%%%%%%%%%%%%%%%%%%%%%%%%%%%%%%%%%%%%%%%%%%%%%%%%%%%%%%%%%%%%%%%%%%%%
% Warm-up exercises (end-of-section, immediate answer)
%%%%%%%%%%%%%%%%%%%%%%%%%%%%%%%%%%%%%%%%%%%%%%%%%%%%%%%%%%%%%%%%%%%%%%%%%%%%%%%

\newcounter{warmup}[chapter]
\renewcommand{\thewarmup}{\thechapter.\arabic{warmup}}

\newenvironment{warmup}[1][\coffee{1}]{%
  \refstepcounter{warmup}%
  \par\medskip\noindent%
  {\sffamily\bfseries Warm-up \thewarmup\ \,#1.}\quad%
  \itshape\ignorespaces%
}{%
  \par\medskip%
}

\newenvironment{warmupSolution}{%
  \par\smallskip\noindent%
  {\sffamily\bfseries\small\itshape Answer:}\quad%
  \small\ignorespaces%
}{%
  \par\smallskip%
}

%%%%%%%%%%%%%%%%%%%%%%%%%%%%%%%%%%%%%%%%%%%%%%%%%%%%%%%%%%%%%%%%%%%%%%%%%%%%%%%
% End-of-chapter exercises (hints in margin, full solution in Appendix A)
%%%%%%%%%%%%%%%%%%%%%%%%%%%%%%%%%%%%%%%%%%%%%%%%%%%%%%%%%%%%%%%%%%%%%%%%%%%%%%%

\newcounter{exercise}[chapter]
\renewcommand{\theexercise}{\thechapter.\arabic{exercise}}

\newenvironment{exercise}[1][\coffee{1}]{%
  \refstepcounter{exercise}%
  \par\medskip\noindent%
  {\sffamily\bfseries Exercise \theexercise\ \,#1.}\quad%
  \ignorespaces%
}{%
  \par\medskip%
}

\newcommand{\hint}[1]{\marginnote{\sffamily\footnotesize\textit{Hint.} #1}}

%%%%%%%%%%%%%%%%%%%%%%%%%%%%%%%%%%%%%%%%%%%%%%%%%%%%%%%%%%%%%%%%%%%%%%%%%%%%%%%
% QED symbol
%%%%%%%%%%%%%%%%%%%%%%%%%%%%%%%%%%%%%%%%%%%%%%%%%%%%%%%%%%%%%%%%%%%%%%%%%%%%%%%

\renewcommand{\qedsymbol}{\ensuremath{\blacksquare}}
